% Generated draft: requires mathematical and source review.
\documentclass[11pt]{article}
\usepackage[margin=25mm]{geometry}
\usepackage{fontspec}
\setmainfont{Latin Modern Roman}
\usepackage{amsmath,amssymb,mathtools,unicode-math}
\setmathfont{Latin Modern Math}
\usepackage{enumitem,needspace,xurl,hyperref}
\setlength{\parindent}{0pt}
\setlength{\parskip}{5pt}
\setlength{\emergencystretch}{3em}
\newcommand{\N}{\mathbb{N}}
\newcommand{\Z}{\mathbb{Z}}
\newcommand{\Q}{\mathbb{Q}}
\newcommand{\R}{\mathbb{R}}
\newcommand{\C}{\mathbb{C}}
\newcommand{\F}{\mathbb{F}}
\newcommand{\A}{\mathbb{A}}
\newcommand{\PP}{\mathbb{P}}
\newcommand{\E}{\mathbb{E}}
\newcommand{\eps}{\varepsilon}
\newcommand{\dd}{\,\mathrm{d}}
\begin{document}
\section*{UnsolvedMath ID: 20002841}
\textbf{Source problem code:} catalogue entry
\subsection{Definitions and mathematical statement}
\textbf{Approximate Conjecture 6, with all appended discussion retained.} Fix a residual representation \(\bar\rho:G_\Q\to\operatorname{GL}_n(k)\), \(k\) a splitting field of characteristic \(p\), satisfying the intended global hypotheses, which are not fully reproduced. For a finite-dimensional coefficient module \(\sigma\) of \(\operatorname{GL}_n(\F_p)\), the source predicts an attached Hecke eigenclass in some degree \(d\) of \(H^d(\Gamma_1(N),\sigma)\), with \(p\nmid N\), exactly when
\[
 N(\bar\rho)\mid N,\qquad
 \operatorname{JH}(\sigma)\cap\operatorname{Wt}(\bar\rho|_{I_p})\ne\varnothing.
\]
Here \(N(\bar\rho)\) is the intended prime-to-\(p\) conductor, \(\operatorname{JH}\) denotes simple constituents, and \(I_p\) is inertia. The arithmetic subgroup, global parity/irreducibility conditions, weight recipe and Hecke normalization are not specified. The source explicitly labels the conjecture approximate and does not choose a single cohomological degree; it suggests a middle range without defining it.

\textbf{Reported attachment and filtration repairs.} In rank at least three, the aid requires attachment through the full unramified Hecke polynomial, with normalized eigenvalues giving every coefficient of
\[
 \det(T-\bar\rho(\operatorname{Frob}_\ell)),\qquad\ell\nmid Np,
\]
rather than only the trace coefficient \(T_\ell\). Even trace and determinant together do not determine a rank-three characteristic polynomial: over an algebraic closure, the invertible diagonal matrices with characteristic polynomials \((T-1)^3\) and \(T^3-3T^2+cT-1\), \(c\ne3\), have the same trace and determinant but different middle coefficient. This is a pointwise matrix diagnostic, not an exhibited pair of global representations with identical data at all primes. Exact powers of \(\ell\) in the Hecke normalization remain unstated.

For a coefficient short exact sequence, the Hecke-equivariant long exact cohomology sequence implies one direction: eigensupport occurring in the whole module occurs in a Jordan--H\"older factor. The reverse is not formal, because connecting maps may kill a factor's class in the extension. Thus the displayed arbitrary-\(\sigma\) equivalence needs additional input beyond the irreducible-weight statement. It is not repaired merely by listing the factors.

\textbf{Preserved Remark 3 and characteristic-zero comparison.} The source hopes that
\[
 \operatorname{Wt}(\bar\rho|_{I_p})
 \subseteq\operatorname{Wt}(\bar\rho|_{I_p}^{\mathrm{ss}}),
\]
mentions an Ash et al. rank-three recipe, and says Florian generally considers the semisimple case; no complete references/definitions are supplied. It asks for analogous assertions over an imaginary quadratic \(E\) using a unitary group that becomes an inner form of \(\operatorname{GL}_n\) over \(E\) and has \(G(\Q_p)\cong\operatorname{GL}_n(\Q_p)\). The Galois representation is then \(\bar\rho:G_E\to\operatorname{GL}_n(k)\), with polarization \(\bar\rho^c\cong\bar\rho^\vee\otimes\nu\) for an unspecified character \(\nu\). The source reports constructing Galois representations in that setting and asks whether either direction of the conjecture can be proved; its exact geometric/global hypotheses are absent.

For dominant integral \(a_1\geq\cdots\geq a_n\), let \(W_a\) be the characteristic-zero highest-weight representation. The source expects a corresponding characteristic-zero eigenclass to have a continuous de Rham Galois representation, with matching Hecke data and Hodge--Tate numbers
\[
 a_1+n-1,\ a_2+n-2,\ldots,\ a_n
\]
in its sign convention. They are distinct. It writes \(\Q_p\) coefficients, but an eigenvalue field may require an extension. For \(n=2\), its example is \(W_a=\operatorname{Sym}^{a_1-a_2}(\mathrm{Std})\otimes\det^{a_2}\). Reduction of a chosen integral lattice need not remain irreducible. Its statement of a unique simple submodule uses an unstated Weyl-versus-dual-Weyl convention: a Weyl module ordinarily has a unique simple head, while the dual convention has a simple socle. No inconsistent lattice identification is silently adopted.

\textbf{Finite-field labels and the reported central-twist audit.} For \(q=p^f\), simple algebraic highest-weight modules \(F_a\) restrict to the usual finite-group labels with \(0\leq a_i-a_{i+1}<q\), modulo simultaneous central shifts by \((q-1)(1,\ldots,1)\). The source's regular subset additionally requires all base-\(p\) digits of each difference to be less than \(p-1\). Put \(C_f=(q-1)/(p-1)\). The printed coordinate formula is
\[
 R(a)_i=a_{n+1-i}-(n-i)C_f.
\]
The aid reports the exact calculation
\[
 R^2(a)=a-(n-1)C_f(1,\ldots,1).
\]
Thus it is generally a square equal to a central determinant twist, not a literal involution. The twist is trivial modulo \(q-1\) only when the corresponding divisibility holds. Selecting dominant regular representatives after the source's coordinate congruences requires extra conventions; the source warns that a regular weight is not always assigned. For \(q=p\), the formula is \((a_n-(n-1),\ldots,a_1)\). The corrupted rank-three diagram labels top \((p-1,p-1)\), bottom \((0,0)\), irregular boundary nodes and an inner ``reducible'' region, and describes a reflection. Its missing edges/dashed line and the precise module meant by reducible are not reconstructed from OCR art; the central-twist formula qualifies the reflection language.

\textbf{Preserved semisimple inertia recipe.} The notes describe a tame semisimple decomposition
\[
 \bar\rho|_{I_p}^{\mathrm{ss}}=
 \bigoplus_j(\chi_j\oplus\chi_j^p\oplus\cdots\oplus\chi_j^{p^{s_j-1}}),
 \qquad\sum_js_j=n,
\]
where \(\chi_j\) factors through \(\F_{p^{s_j}}^*\) and \(s_j\) is its minimal Frobenius orbit length. They form the torus \(T\cong\prod_j\F_{p^{s_j}}^*\) in \(\operatorname{GL}_n(\F_p)\), lift characters to characteristic zero, and attach \(V(T,\widetilde\chi)\). Regularity, signs and degeneracies needed for this Deligne--Lusztig/cuspidal-induction description are not fully given, so ``usually irreducible'' is not an unconditional assertion. In the rank-three split case the indicated principal series is induction from the torus product character \(\widetilde\chi_1\otimes\widetilde\chi_2\otimes\widetilde\chi_3\) inflated to the Borel, not a direct sum of Borel characters. In the niveau-three case the source gives, on \(\alpha\in\F_{p^3}^*\setminus\F_p^*\), the character value
\[
 \widetilde\chi(\alpha)+\widetilde\chi(\alpha)^p+
 \widetilde\chi(\alpha)^{p^2},
\]
under its chosen cuspidal/character convention. This is part of an incomplete recipe, not a full definition of \(\operatorname{Wt}\).

The aid proves only these normalization/filtration audits and leaves the global modularity-and-weight equivalence unresolved. The long appended remarks, historical assertions and reconstructed formulas require independent mathematical and source verification.
\subsection{Short English statement}
The aid audits an approximate higher-rank Serre-weight conjecture: use the full Hecke polynomial, keep the Jordan--H\"older implication one-way without extra input, and account for the central twist in the printed weight map. All lengthy unitary, highest-weight and inertia-recipe remarks are retained with their missing conventions, without claiming the global conjecture.
\subsection{Sources}
\begin{itemize}
\item Local UnsolvedMath registry, permanent ID \(20002841\), catalogue entry, \texttt{aim-representation-theory-notes.json}, \texttt{notes[112]}, Michael Volpato's \emph{\(p\)-adic representations, modularity, and beyond} notes, pp. 21--24, approximate Conjecture 6 and Remark 3: \url{https://aimath.org/WWN/padicmodularity/padicmodularity.pdf}.
\item Original discussion mentions Ash et al., Florian, unitary groups, highest-weight modules and finite-field torus/cuspidal recipes without full primary citations. The diagram and several formulas have substantial extraction damage.
\item Supplied research summary on full-Hecke attachment, the filtration obstruction and the \(q=p^f\) central-twist calculation. Attempt 1: partial\_result, status at run ``open'', rigor ``proved'', low novelty confidence; registry partially\_solved and reconstructed\_unverified. The machine-generated aid is not peer reviewed and must be independently verified before citation; no PDF visit is claimed.
\end{itemize}
\end{document}
